VSM is a session-level benchmark for Vietnamese student mental-health support. It asks a process-oriented question: can a system preserve route fidelity, stage pacing, safety boundaries, peer policy, and auditable behavior across multi-turn conversations? The main split, \texttt{vsm\_session\_core.jsonl}, contains 100 sessions and 1200 user turns. \texttt{vsm\_probe.jsonl} is reserved for engineering checks, and \texttt{vsm\_stress.jsonl} is reserved for robustness analysis.

\paragraph{Cases and tasks.}
Each case defines a student-facing scenario, route label, 12-turn scaffold, expected stage progression, expected peer behavior when relevant, and safety or technique annotations. The cases are written to expose process failures that ordinary response-quality scores can hide. CBT cases test validation, ABC capture, distortion labeling, Socratic reasoning, and action planning. Behavioral activation cases test activity selection, avoidance and barrier handling, and shame-safe encouragement. Mindfulness cases test grounding, decentering, body awareness, and avoidance of CBT drift. Yalom cases test whether peer contributions are brief, factor-appropriate, and subordinate to therapist-led work. Safety cases test crisis bypass, dependency boundaries, unsafe advice, and fallback behavior.

\begin{table}[t]
\centering
\footnotesize
\setlength{\tabcolsep}{3pt}
\begin{tabular}{p{0.2\linewidth}p{0.31\linewidth}p{0.35\linewidth}}
\toprule
Route & Main task & Primary checks \\
\midrule
CBT & Pace thoughts, distortions, Socratic reasoning, and action. & stage accuracy, technique fidelity, peer silence/use \\
BA & Move from energy check to micro-action and barrier planning. & actionability, activation fit, shame-safe language \\
MBI & Support grounding, decentering, body awareness, and mindful action. & route fidelity, non-CBT drift, embodied guidance \\
Yalom & Test peer support factors such as universality and hope. & peer selection, persona validity, repetition \\
Safety & Handle self-harm, dependency, and boundary-sensitive inputs. & unsafe advice, crisis bypass, fallback behavior \\
\bottomrule
\end{tabular}
\caption{Planned VSM route/task definition used for the main benchmark table after the full run.}
\label{tab:vsm_route_tasks}
\end{table}


\paragraph{Outputs and metadata.}
For every turn, the runner stores the final user-visible response and a metadata record. \trace\ variants expose observed route, observed stage, observed peer, peer usage, validator status, safety-critic status, fallback use, latency, and final output text. Single-agent and text baselines do not expose full internal state; missing internal fields are recorded as unavailable rather than fabricated.

\paragraph{Scoring.}
VSM separates reliability failures from capability mismatches. Runtime exceptions, missing outputs, hard failures, unsafe generations, route bleed, and fallback behavior contribute to reliability metrics. Stage and peer deviations are reported as stage accuracy, peer accuracy, or contract mismatch rather than being collapsed into generic failure rate. This distinction matters for ablations: \textsc{TRACE-NoPeer} should lose peer-dynamics credit, but it should not be counted as a runtime failure merely because it intentionally omits peer agents.

\begin{table*}[t]
\centering
\footnotesize
\setlength{\tabcolsep}{4pt}
\begin{tabular}{p{0.17\linewidth}p{0.24\linewidth}p{0.32\linewidth}p{0.17\linewidth}}
\toprule
Metric group & Source & Definition & Reported as \\
\midrule
Reliability & Runner logs & Completed turns, missing outputs, runtime exceptions, hard failures, unsafe generations, route bleed, and fallback use. & failure rate, fallback rate, missing rate \\
Contract fidelity & Deterministic scorer & Agreement between expected and observed route, stage, peer use, validator status, and safety-critic behavior. & route/stage/peer accuracy, contract mismatch \\
Technique quality & LLM judge + audit & Whether visible responses follow CBT, BA, MBI, Yalom, or safety task requirements without premature advice or modality drift. & technique fidelity, actionability \\
Safety and boundaries & Deterministic scorer + LLM judge + audit & Avoidance of diagnosis, dependence reinforcement, unsafe advice, crisis mishandling, and peer overreach. & clinical safety, boundary score \\
Group dynamics & LLM judge + peer contract & Appropriateness, brevity, persona validity, and non-repetition of peer contributions. & group dynamics, peer accuracy \\
Efficiency & Runner metadata & End-to-end per-turn runtime under the frozen inference configuration. & average latency, p95 latency \\
\bottomrule
\end{tabular}
\caption{VSM scoring views. Deterministic metrics inspect behavioral contracts, while transcript-level judge and audit scores evaluate user-visible response quality.}
\label{tab:vsm_metrics}
\end{table*}


\paragraph{Benchmark status.}
The paper treats VSM as a frozen evaluation protocol until the full core run is complete. Result tables in Section~\ref{sec:result_placeholders} are intentionally schema-only slots: they define what will be reported, but they do not contain evidence from partial or engineering-only runs.
